\subsection{Masas, hijo estacionario y restos cúbicos}\label{sec:B5T03}
Escribimos $|X|_\Sigma=\sum_{ij}|X_{ij}|$ para matrices y
$|x|_1=\sum_i|x_i|$ para vectores. Son controles de masa, no de dimensión.
Las identidades de incidencia son
\[
 \begin{gathered}
 \sum_ep_e=k,\qquad\sum_d m_d=2k,\qquad
 m_{r,a}=\sum_{e\ni a}p_e^r,\\
 \sum_a m_{2,a}=4\sum_ep_e^2=O(1),\qquad
 \max_a m_{2,a}=O(k^{-1}).
 \end{gathered}
\]
Al insertar $c$ se retiran sus cuatro líneas, $J$, y las casillas de su unión, $F$.
Cada línea retenida pierde a lo sumo tres casillas. Por tanto
\[
 |F|\le4k-3,\quad\sum_Fp=O(1),\quad\sum_Fp^3=O(k^{-2}),
 \quad d_a\le3C/k,\quad |d|_1=O(1).
\]

\paragraph{Reequilibrado y paso de Newton.}
El reequilibrado RC de la Sección~\ref{sec:B3T07} usa únicamente $2k$ líneas: al restringir
el padre, la brecha pierde a lo sumo $6C/k$. Resolver su laplaciano da
\[
 w=p|_{E'},\quad\bar p=w+z,\quad B\bar p=\mathbf1,
 \quad z/w\in\operatorname{im}B^T,\quad\|z/w\|_\infty=O(k^{-1}).
\]
Aquí $B$ es la incidencia RC y $D$ la incidencia diagonal del hijo.
Con $u=d-D^Tz$ se obtiene
\[
 |z|_1=O(1),\quad\|u\|_\infty=O(k^{-1}),\quad|u|_1=O(1),
 \quad\sum u^2=O(k^{-1}),\quad\sum|u|^3=O(k^{-2}).
\]
Las vacancias de todo el segmento relevante son al menos $\eta/2$.
El residuo KKT, módulo funciones RC, es
\[
 e=D(u/h)+O_\infty(k^{-2})=O_\infty(k^{-1}).
\]
La Sección~\ref{sec:B5T02} da una respuesta tangente $\bar{\mathcal R}$ con norma relativa acotada.
Para $\|v/\bar p\|_\infty\le r$,
\[
 \nabla H(\bar p+v)=\nabla H(\bar p)-Kv+\mathcal N(v),\quad
 \|\mathcal N(v)\|_\infty\le Cr^2,
\]
y su constante de Lipschitz relativa es $Cr$.
La ecuación $v=\bar{\mathcal R}e+\bar{\mathcal R}\mathcal N(v)$
es contractiva en radio $O(k^{-1})$. Este argumento de Newton local da un hijo estacionario único
en la rama pequeña:
\[
 \|(p'-\bar p)/\bar p\|_\infty=O(k^{-1}),\quad
 \|(p'-\bar p-\bar{\mathcal R}e)/\bar p\|_\infty=O(k^{-2}).
\]
En particular $y=p'-w$ cumple $y/p=O(k^{-1})$ y $|y|_1=O(1)$.
La herencia es de una clase interior a una exterior con margen;
iterar esta afirmación no demuestra permanencia en una clase fija.

\paragraph{Restos cúbicos sin contar las diagonales.}
Si $|v_e|\le rp_e$, entonces
\[
 \sum_d\frac{|(D^Tv)_d|^3}{h_d^2}
 \le\eta^{-2}r^3\sum_d m_d^3
 \le2\eta^{-2}\varepsilon^2kr^3.
\]
La parte de casillas satisface $\sum|v|^3/p^2\le kr^3$.
Para $r=O(k^{-1})$, ambos restos son $O(k^{-2})$.
Las casillas retiradas no entran en este desarrollo de Taylor: para ellas se usa la identidad
de frontera exacta de la Sección~\ref{sec:B4T04}.

\paragraph{Diferencia de Schur completa.}
Con $Q=S_{IJ}S_{JJ}^{-1}S_{JI}$ y $U=A_{F,I}^TT_FA_{F,I}$,
la diferencia exacta es $\mathcal E=Q+\operatorname{diag}_D d-U+\Xi$,
donde $\Xi$ corresponde al reequilibrado $z$.
Cada columna de $S_{IJ}$ tiene masa $O(1)$ y sólo hay cuatro columnas.
Se sigue que $|\mathcal E|_\Sigma=O(1)$ y $\|\mathcal E\|_\infty=O(k^{-1})$,
por lo que $\|\mathcal E\|_{HS}^2=O(k^{-1})$.
Además $\sum_a|\mathcal E_{aa}|=O(k^{-1})$:
la diagonal de un punto factible es $1+3m_{2,a}+O(m_{3,a})$,
$\sum_a|\Delta m_{2,a}|=O(k^{-1})$ y $\operatorname{tr}Q=O(k^{-1})$.
Con ello, la estimación de la traza lineal tampoco depende de la dimensión.
